\documentclass[instructions.tex]{subfiles}
\begin{document}
 
\chapter{On doit parler avec respect du Pape et des Évêques.}
 
\primo Il y a peu de gens qui aiment cet esprit de droiture que le Prophète demandait à Dieu. Combien de gens font les beaux esprits, en blâmant tout, en critiquant même ce qu'il y a de plus respectable dans l'Église! On dirait que les premiers pasteurs ne sont sur le chandelier que pour être en butte à la contradiction et à la censure. On veut raisonner du gouvernement de l'Église et du souverain pontife; et comment en parle-t-on? sans respect, sans charité, sans discrétion; témérité qui ne peut manquer d'avoir de fâcheuses suites pour la religion. Dieu veuille nous préserver de voir parmi nous ce que le prophète Isaïe a prédit : Ces hommes imprudens, sans sagesse, qui, par des raisonnemens pleins de fierté, séduisent ceux qui ne comprennent pas l'artifice de leur langage.
 
A force d'entendre parler désavantageusement du pape, on s'y accoutume; on perd peu à peu le respect et la confiance qu'on doit avoir pour le père commun des fidèles; on quitte insensiblement l'esprit de subordination, sans lequel la religion ne peut subsister. Luther ne trouva point de moyen plus efficace pour répandre son hérésie et autoriser son schisme, que de se déchaîner contre le pape et les prétendus abus de l'Église romaine.
 
Tel a été dans tous les temps l'esprit des ennemis de l'Église : ils ont supposé des faits calomnieux, altéré l'histoire ecclésiastique, répandu dans leurs livres des réflexions pleines de malignité contre la conduite des papes et l'autorité du Saint-Siége. C'est ainsi qu'en rendant odieux les souverains pontifes, ils ont séparé les fidèles du chef de l'Église, et divisé le troupeau de Jésus-Christ. Quand on s'en prend au pasteur, les ouailles sont bientôt dispersées\footnote{Percutiam pastorem, et dispergentur oves. Marc. 14. 25}.
 
\secundo Il n'est pas moins dangereux de parler mal des évêques et des autres supérieurs ecclésiastiques. Se donner la liberté de parler avec si peu de ménagement du gouvernement de l'Église, de la conduite et des lois de ses supérieurs, c'est marquer peu de religion dans le coeur, et beaucoup de légèreté dans l'esprit. De quoi vous mêlez-vous, disait-on à un philosophe, de censurer le gouvernement et les lois que les magistrats établissent pour le bon ordre, vous qui ne pouvez établir l'ordre dans votre maison, ni mettre en paix votre servante et votre femme? Tel qui ne sait pas seulement gouverner sa famille, ni régler sa maison, se mêle de critiquer le gouvernement ecclésiastique et les ordonnances de son prélat. Combien de gens, qui blâment dans un supérieur ou dans un évêque, ce qui fait son mérite et son éloge!
 
Après tout, quand nous apercevrions quelque trait répréhensible dans un supérieur, cela autoriserait-il notre censure et notre indocilité? Un roi est-il moins notre maître, est-il moins en droit de commander, parce qu'il n'est pas saint? De même les évêques sont-ils moins supérieurs, parce qu'ils sont hommes? et leur devons-nous moins l'obéissance, parce qu'ils ne sont point impeccables?
 
Si l'on voyait un ouvrier sur une tour fort élevée, aussitôt qu'on le verrait chanceler, n'en aurait-on pas compassion? s'aviserait-on de le railler, parce qu'il ferait quelques faux pas? Le fardeau et l'élévation des premiers pasteurs de l'Église, loin d'exciter nos médisances et nos satires, devraient bien plutôt exciter notre charité et nos prières. Craignons Dieu, dit saint Bernard, et prions pour ceux dont l'élévation et le ministère font trembler les anges mêmes.
 
\section{Résolutions}
 
\primo Je regarderai toujours le pape comme chef de l'Église, auquel les fidèles doivent toute obéissance comme à leur père commun.
 
\secundo Et je respecterai les évêques comme successeurs des apôtres, auxquels Jésus-Christ a dit : Celui qui vous écoute m'écoute; celui qui vous méprise me méprise... 

\prayer{Mon Dieu, vous avez promis d'être tous les jours avec le corps des pasteurs de votre Église; nous comptons avec confiance sur votre infaillible promesse.}
 
\end{document}
